\documentclass[10pt]{article}
\usepackage[margin=0.7in]{geometry}
\usepackage{xcolor}
\usepackage[hidelinks]{hyperref}

\definecolor{spark}{HTML}{C2410C}
\setlength{\parindent}{0pt}
\setlength{\parskip}{3pt}
\pagestyle{empty}

% compact section heading (no titlesec dependency)
\newcommand{\sechead}[1]{\vspace{3pt}\par{\color{spark}\large\bfseries #1}\par\vspace{1pt}}
% compact bullet list (no enumitem dependency)
\newenvironment{tight}{\begin{itemize}\setlength{\itemsep}{1pt}\setlength{\parskip}{0pt}\setlength{\topsep}{1pt}\setlength{\leftmargin}{1.2em}}{\end{itemize}}
\newcommand{\lead}[1]{\textbf{\color{spark}#1}\;}

\begin{document}

{\LARGE\bfseries Spark Estimator}\hfill{\large Ashrit Kommireddy}\\[1pt]
{\color{spark}\rule{\textwidth}{1.2pt}}\\[2pt]
Repair Cost Estimator --- Project Writeup \hfill
Live: \href{https://spark-estimator.vercel.app/}{spark-estimator.vercel.app}\\
\hspace*{\fill}Source: \href{https://github.com/ashritkvs/spark-estimator}{github.com/ashritkvs/spark-estimator}

\sechead{Most interesting design decision}
I optimized for a real walkthrough, not a screenshot. An agent stands in an empty house and
types quantities fast, so the running total had to update with \emph{zero} jank. Rather than
re-render the whole screen on every keystroke (which loses input focus and stutters), I built a
two-path render: structural changes (checking an item, expanding a group, switching rooms) do a
simple full re-render, but quantity entry \emph{patches only the affected DOM nodes} --- the line
total, its group total, the header total, and the progress bar. The result feels native while the
code stays simple. The same ``field-first'' instinct led me to drop the reference app's Tailwind
play-CDN and hand-write the CSS, so the app paints instantly and is genuinely offline (no 3\,MB
runtime compiler to cache). Internally, everything keys off the price-list \texttt{id} and a
\texttt{room::item} state model, so the CSV stays the single source of truth and rooms can be
cloned freely.

\sechead{What is broken or fragile}
\begin{tight}
  \item \lead{Serial OCR.} The age decoder is an honest best-effort heuristic across common
  HVAC / appliance serial formats (explicit year, Carrier week-year, Rheem year-week, 2-digit
  prefixes), not an exhaustive per-brand database --- so it always lets the agent confirm or
  override the detected year. Tesseract also needs one online load before it can run offline.
  \item \lead{Photo storage.} Photos live in \texttt{localStorage} (downscaled to 1280\,px). A very
  photo-heavy project can approach the browser quota; the app degrades by keeping selections and
  dropping the heaviest photo data rather than failing the save. A production build would move
  photos to IndexedDB.
  \item \lead{Local-only.} State is per-device. There is no multi-agent sync yet.
\end{tight}

\sechead{Creative additions (and why)}
My required creative addition --- a feature not listed in the brief --- is the
\textbf{Deal Analyzer}. I also extended the required photo capture into a serial age decoder.
Both target the part of the job the brief says hurts most: missed hidden costs and mispriced offers.
\begin{tight}
  \item \lead{Deal Analyzer.} The live repair total flows into a Maximum Allowable Offer
  ($\mathrm{MAO} = \mathrm{ARV}\times(1-\mathrm{margin}) - \mathrm{repairs} - \mathrm{costs}$), projected
  profit, and a GO / TIGHT / NO verdict. It connects every checkbox in the walkthrough to the only
  number that decides the deal, and rides along into the Excel export. I chose it because the brief
  frames the whole business around ARV minus repairs minus offer --- this makes that math live.
  \item \lead{Serial-photo OCR + equipment age decoder.} Scan an HVAC / water-heater / appliance
  nameplate; the app reads the serial, decodes the manufacture year, and flags units past their
  service life --- catching the missed furnace or water heater that swings an estimate by thousands.
\end{tight}

\sechead{What I would ship next (two more days)}
\begin{tight}
  \item Cloud sync so a team shares projects across devices, with offline-first conflict handling.
  \item A materials / scope-of-work export grouped by trade, generated from the same selections.
  \item ARV assist: pull a rough comp-based ARV so the Deal Analyzer needs less manual input.
  \item Harden OCR with a per-brand decode library and an on-device confidence score.
\end{tight}

\sechead{Role of AI tools}
I built this with Claude Code (Opus) as a pair programmer. I owned the architecture, the data model,
the two creative features, the offer math, and all product decisions; AI accelerated scaffolding the
107-item data set from the CSV, iterating on the mobile UI, and debugging the export and
service-worker caching. Every line was reviewed by me and the full flow was verified in-browser
(live totals, ZIP/Excel export, OCR, offline) before shipping.

\end{document}
